\chapter*{Conclusions and Future Work}
\label{cap:conclusions}
\addcontentsline{toc}{chapter}{Conclusions and Future Work}

The objective of this study was to create a tool that, given a set of tasks and conditions, would be capable of generating behavior trees for non-player characters (NPCs) to serve as a foundation to reduce the workload of video game designers.

Based on the aforementioned results, the following conclusions have been reached.

\section{Conclusions}

This work demonstrates the viability of generating functional behavior trees (BTs) for NPCs in Unreal Engine 5 from natural language instructions.
Through a hybrid architecture that combines Mistral Large 3 for the logical structure and llama3:8B for Blackboard variable assignment, the system successfully mitigates hallucinations and translates the generated pseudocode into readable JSON formats.

To ensure behavior quality, a double automatic validation system was implemented (a nondeterministic finite automaton (NFA) to check dependencies and a final objective validator).
Experimental tests with designers (N = 3) confirmed the tool's high usability, achieving an excellent System Usability Scale (SUS) score of 90.83/100 and demonstrating a drastic reduction in development times compared to manual creation.

Although the automatically generated trees exhibit lower qualitative quality due to the inclusion of some redundant nodes, their execution in the game meets user expectations.
The tool establishes itself as an excellent starting point that facilitates subsequent editing and significantly decreases the designer's workload.
Despite the quantitative results being exploratory due to the sample size, the project successfully validates the integration of LLMs into game production workflows.

\section{Future Work}

Given the limitations presented in section \ref{sec:Limitaciones}, the future work of this study can focus on several efforts:

\begin{itemize}
    \item With greater computational capabilities, study the use of larger local models for the tasks performed, as well as applying fine-tuning techniques to verify if their use would match or improve current results in order to avoid dependency on third-party APIs.
    \item Improve the RAG system by incorporating video game BTs, saving the vector store on disk, and enabling metadata searches using Qdrant, which could potentially improve the context provided to the models and reduce latency.
    \item Implement an LLM agent with a visual interface to modify the generated BTs more easily, allowing natural language feedback to be provided to the generator.
    \item Implement an LLM agent capable of reading the project's actions and descriptions to automate their subscription to the Abstract Factory and generate descriptions for the tool.
    \item Implement an LLM agent to automate the validation tests by generating the dependency order and a goal function.
    \item Gather more users for the trials in order to achieve a statistically significant sample to definitively conclude the proposed hypotheses.
\end{itemize}


